\section{Conclusion}\label{sec:conclusion}

We measured how much of a person's own transaction history an unsupervised behavioural anomaly detector needs, on real bank histories, in a controlled simulation and on a structurally different synthetic card dataset, using a ground-truth protocol defined by provenance rather than by any statistical rule. A sufficiency point exists and arrives earlier than the framing of the question suggests: about 25 personal transactions for ranking quality with the best real-data detector, saturation by 50, and the same shape, with a later threshold, on card-structured data whose global baseline is much stronger. Below that point a population model built from other people's behaviour does better than the person's own history, and a shrinkage hybrid that degrades to the population when history is thin removes the penalty at no measurable cost later.

The single number is the least useful part of the answer. Whether a ranking-quality threshold translates into a usable alert threshold depends on how the alert threshold is calibrated: scoring a model's own training data makes it flag 96\% of transactions at 10 personal transactions, and leave-one-out calibration makes the same model nearly silent there and near-nominal from 25. Detector rankings under scarcity are dataset-dependent, behavioural variability changes the requirement tenfold when manipulated but is not detectable across real accounts, and a direct drift measurement does not explain why older history hurts. No history size made performance stable across individuals, which undermines any per-person guarantee derived from a population average.

Two extensions would sharpen these results. A finer history grid between 10 and 25, together with alternative out-of-sample threshold rules (split-half and population-derived) on both datasets, would localise the sufficiency point and test whether LOO calibration is the right rule. A targeted evaluation on entities with historical generator-labelled fraud, and a deeper real dataset carrying merchant fields and individual-level demographics, would test the findings against fraud labels and enable an individual-level fairness audit.
